% =====================================================================
%  PRÉAMBULE COMMUN — Carnet d'entraînement, Fiche 2 (Nomenclature)
%  (inséré physiquement dans chaque document pour qu'il soit autonome)
% =====================================================================
\documentclass[11pt,a4paper]{article}

% ---------- Encodage & langue ----------
\usepackage[utf8]{inputenc}
\usepackage[T1]{fontenc}
\usepackage{lmodern}
\usepackage[french]{babel}

% ---------- Mathématiques & chimie ----------
\usepackage{amsmath,amssymb,mathtools}
\usepackage{icomma}
\usepackage[version=4]{mhchem}
\usepackage{siunitx}
\sisetup{output-decimal-marker={,},per-mode=symbol}

% ---------- Graphismes & mise en page ----------
\usepackage{xcolor}
\usepackage{colortbl}
\usepackage{tikz}
\usepackage[most]{tcolorbox}
\usepackage{enumitem}
\usepackage{array}
\usepackage{booktabs}
\usepackage{multicol}
\usepackage{fancyhdr}
\usepackage{titlesec}
\usepackage[hidelinks]{hyperref}
\usetikzlibrary{arrows.meta,positioning,calc,shapes.geometric,shapes.misc,
  fit,backgrounds,decorations.pathreplacing}
\usepackage[a4paper,margin=2.2cm,headheight=28pt,headsep=10pt]{geometry}

% =====================================================================
%  PALETTE (charte « Chimie » — mauve/prune du cours)
% =====================================================================
\definecolor{primary}{RGB}{146,95,161}
\definecolor{primarydark}{RGB}{92,55,108}
\definecolor{defcol}{RGB}{0,114,178}
\definecolor{thmcol}{RGB}{0,140,120}
\definecolor{formulacol}{RGB}{198,93,9}
\definecolor{remarkcol}{RGB}{120,94,200}
\definecolor{attncol}{RGB}{200,40,55}
\definecolor{excol}{RGB}{0,135,90}
\definecolor{methcol}{RGB}{90,90,100}
\definecolor{metalcol}{RGB}{198,150,40}
\definecolor{nonmetcol}{RGB}{0,140,150}
\definecolor{oxycol}{RGB}{210,55,45}
\definecolor{hydcol}{RGB}{60,120,210}
\definecolor{catcol}{RGB}{200,70,120}
\definecolor{ancol}{RGB}{40,110,190}
\definecolor{saltcol}{RGB}{120,90,160}

% =====================================================================
%  STYLE COMMUN DES BOÎTES
% =====================================================================
\tcbset{
  boxbase/.style={enhanced,breakable,boxrule=0.9pt,arc=2.5pt,
     left=3mm,right=3mm,top=2.3mm,bottom=2.3mm,
     fonttitle=\bfseries,coltitle=white,
     toptitle=0.6mm,bottomtitle=0.6mm},
}

\newtcolorbox{definition}[1][]{%
  boxbase,colback=defcol!5,colframe=defcol,
  title={\textsc{Définition}\ifx\relax#1\relax\relax\else\textmd{ --- \itshape #1}\fi}}
\newtcolorbox{regle}[1][]{%
  boxbase,colback=thmcol!6,colframe=thmcol,
  title={\textsc{Règle}\ifx\relax#1\relax\relax\else\textmd{ --- \itshape #1}\fi}}
\newtcolorbox{formule}[1][]{%
  boxbase,colback=formulacol!6,colframe=formulacol,
  title={\textsc{Méthode-clé}\ifx\relax#1\relax\relax\else\textmd{ --- \itshape #1}\fi}}
\newtcolorbox{remarque}[1][]{%
  boxbase,colback=remarkcol!6,colframe=remarkcol,
  title={\textsc{Remarque}\ifx\relax#1\relax\relax\else\textmd{ : \itshape #1}\fi}}
\newtcolorbox{attention}[1][]{%
  boxbase,colback=attncol!6,colframe=attncol,
  title={\textsc{Attention}\ifx\relax#1\relax\relax\else\textmd{ : \itshape #1}\fi}}
\newtcolorbox{exemple}[1][]{%
  boxbase,colback=excol!5,colframe=excol,
  title={\textsc{Exemple}\ifx\relax#1\relax\relax\else\textmd{ : \itshape #1}\fi}}
\newtcolorbox{methode}[1][]{%
  boxbase,colback=methcol!7,colframe=methcol,sharp corners,
  title={\textsc{Méthode}\ifx\relax#1\relax\relax\else\textmd{ : \itshape #1}\fi}}
\newtcolorbox{astuce}[1][]{%
  boxbase,colback=primary!6,colframe=primary,
  title={\textsc{Astuce concours}\ifx\relax#1\relax\relax\else\textmd{ : \itshape #1}\fi}}

% ---- Exercice (numéroté) ----
\newtcolorbox[auto counter]{exercice}[1][]{%
  boxbase,colback=primary!4,colframe=primarydark,
  title={\textsc{Exercice}~\thetcbcounter\ifx\relax#1\relax\relax\else\textmd{ --- \itshape #1}\fi}}

% ---- Corrigé (numéroté) ----
\newtcolorbox[auto counter]{corrige}[1][]{%
  boxbase,colback=excol!4,colframe=excol,
  title={\textsc{Corrigé}~\thetcbcounter\ifx\relax#1\relax\relax\else\textmd{ --- \itshape #1}\fi}}

% =====================================================================
%  EN-TÊTES ET PIEDS DE PAGE
% =====================================================================
\pagestyle{fancy}
\fancyhf{}
\fancyhead[L]{\small\textcolor{primarydark}{\textbf{Fiche 2} — Nomenclature}}
\fancyhead[R]{\small\textcolor{primarydark}{\textsc{Carnet d'entraînement}}}
\fancyfoot[C]{\small\thepage}
\renewcommand{\headrulewidth}{0.8pt}
\renewcommand{\headrule}{\hbox to\headwidth{\color{primary}\leaders\hrule height \headrulewidth\hfill}}

\titleformat{\section}{\Large\bfseries\color{primarydark}}{\thesection}{0.6em}{}
\titleformat{\subsection}{\large\bfseries\color{primary}}{\thesubsection}{0.6em}{}
\titleformat{\subsubsection}{\normalsize\bfseries\color{primary!80!black}}{}{0em}{}

% ---- Bandeau de titre réutilisable : \bandeau{Thème}{Sous-titre} ----
\newcommand{\bandeau}[2]{%
  \begin{center}
  \begin{tikzpicture}
  \node[rectangle,rounded corners=7pt,fill=primary,text=white,
        minimum width=16.0cm,minimum height=1.7cm,align=center,inner sep=8pt]
    {{\Large\bfseries #1}\\[3pt]{\normalsize #2}};
  \end{tikzpicture}
  \end{center}
  \vspace{0.3cm}}
\begin{document}
\bandeau{Thème 1 — Facile}{Ions, valences et nombre d'oxydation --- Corrigé}

\begin{corrige}[métal ou non-métal ?]
Règle : à gauche/centre du tableau $\to$ métal $\to$ cation ; à droite $\to$ non-métal $\to$ anion.
\begin{center}
\begin{tabular}{lcc}
\toprule
Élément & Nature & Ion formé \\
\midrule
\ce{Na} & métal & cation (\ce{Na+}) \\
\ce{Cl} & non-métal & anion (\ce{Cl-}) \\
\ce{Ca} & métal & cation (\ce{Ca^2+}) \\
\ce{O}  & non-métal & anion (\ce{O^2-}) \\
\ce{Al} & métal & cation (\ce{Al^3+}) \\
\ce{S}  & non-métal & anion (\ce{S^2-}) \\
\ce{Fe} & métal & cation (\ce{Fe^2+} ou \ce{Fe^3+}) \\
\ce{F}  & non-métal & anion (\ce{F-}) \\
\bottomrule
\end{tabular}
\end{center}
\end{corrige}

\begin{corrige}[lire une valence]
La valence est la \emph{valeur absolue} de la charge :
\[ \ce{K+}:1,\quad \ce{Mg^2+}:2,\quad \ce{Al^3+}:3,\quad \ce{Br-}:1,\quad \ce{O^2-}:2,\quad \ce{PO4^3-}:3,\quad \ce{NH4+}:1,\quad \ce{SO4^2-}:2. \]
Le signe (+ ou $-$) n'intervient pas dans la valence ; il indique seulement si l'ion est cation ou anion.
\end{corrige}

\begin{corrige}[électrons gagnés ou perdus]
Cation : $\ce{e-} = Z - \text{charge}$. Anion : $\ce{e-} = Z + |\text{charge}|$.
\begin{center}
\begin{tabular}{lccccc}
\toprule
Ion & \ce{Na+} & \ce{Ca^2+} & \ce{Al^3+} & \ce{Cl-} & \ce{S^2-} \\
\midrule
Bilan & perd 1 & perd 2 & perd 3 & gagne 1 & gagne 2 \\
Électrons & $11-1=10$ & $20-2=18$ & $13-3=10$ & $17+1=18$ & $16+2=18$ \\
\bottomrule
\end{tabular}
\end{center}
Remarque : \ce{Na+}, \ce{Al^3+} atteignent 10 électrons (comme le néon) ; \ce{Ca^2+}, \ce{Cl-}, \ce{S^2-}
atteignent 18 (comme l'argon). Les ions cherchent la configuration du gaz noble le plus proche.
\end{corrige}

\begin{corrige}[n.o.\ des cas immédiats]
\begin{enumerate}[label=\alph*),leftmargin=1.8em,topsep=1pt,itemsep=1pt]
  \item \ce{Fe} : corps simple $\Rightarrow$ n.o.\ $= 0$ (convention 1).
  \item \ce{Cl2} : corps simple $\Rightarrow$ n.o.\ $= 0$ (convention 1).
  \item \ce{Ca^2+} : ion monoatomique $\Rightarrow$ n.o.\ $= +\mathrm{II}$ (convention 2).
  \item \ce{Cl-} : ion monoatomique $\Rightarrow$ n.o.\ $= -\mathrm{I}$ (convention 2).
  \item oxygène dans \ce{Na2O} : $-\mathrm{II}$ (convention 3 ; ce n'est pas un peroxyde).
\end{enumerate}
\end{corrige}

\begin{corrige}[repérer le cation]
Le seul \textbf{cation} est \ce{NH4+} (ammonium), car sa charge est \emph{positive}. Tous les autres portent
une charge négative, ce sont des \textbf{anions} : \ce{NO3-} (nitrate), \ce{SO4^2-} (sulfate),
\ce{MnO4-} (permanganate), \ce{CO3^2-} (carbonate), \ce{PO4^3-} (phosphate). \\
En nomenclature, \ce{NH4+} joue le rôle d'un métal de valence 1 (comme \ce{Na+} ou \ce{K+}).
\end{corrige}
\end{document}
